\chapter{Hybrid Text Analysis Pipeline}
\label{app:pii_detector}

In the deployed system, OCR-extracted text was not classified by the NER model
alone. Instead, the final text categorization step was implemented through a
hybrid pipeline in \texttt{pii\_detector.py}, where NER predictions were
combined with rule-based checks, contextual field hints, and post-processing
filters. Listings~\ref{lst:pii_detector_extract} and
\ref{lst:pii_detector_rules} show the central extraction logic and an example of
the rule-based entity detection used in the final system.

\begin{lstlisting}[language=Python, caption={Core hybrid extraction logic used after OCR text recognition}, label={lst:pii_detector_extract}]
def extract_entities(
    text: str,
    expected_label: Optional[str] = None,
) -> List[TextEntity]:
    clean = normalize_text(text)
    if not clean:
        return []

    all_rule_ents: List[TextEntity] = []
    all_gliner_ents: List[TextEntity] = []
    all_custom_ents: List[TextEntity] = []

    all_rule_ents.extend(detect_rule_entities(clean, expected_label))
    all_custom_ents.extend(detect_custom_ner_entities(clean, expected_label))
    all_gliner_ents.extend(detect_gliner_entities(clean, expected_label))

    for line in split_lines(clean):
        key, value = extract_key_value(line)
        line_expected = context_label_from_key(key) or expected_label
        segment = value if key else line

        if not segment or is_noise(segment):
            continue

        all_rule_ents.extend(detect_rule_entities(segment, line_expected))
        all_custom_ents.extend(detect_custom_ner_entities(segment, line_expected))
        all_gliner_ents.extend(detect_gliner_entities(segment, line_expected))

        # ... context-based helper rules omitted ...

    entities = merge_all(all_rule_ents, all_gliner_ents, all_custom_ents)

    if (
        not entities
        and expected_label in {"POSTAL_CODE", "PHONE", "ID_NUMBER", "DATE"}
        and len(clean) >= 2
        and not is_noise(clean)
    ):
        entities = [TextEntity(clean, expected_label, 0.50, "context")]

    return dedupe_entities(entities)
\end{lstlisting}

\begin{lstlisting}[language=Python, caption={Example of rule-based entity detection used alongside NER}, label={lst:pii_detector_rules}]
def detect_rule_entities(
    text: str,
    expected_label: Optional[str] = None,
) -> List[TextEntity]:
    t = normalize_token(text)
    entities: List[TextEntity] = []

    if is_noise(t):
        return entities

    for m in EMAIL_RE.finditer(t):
        entities.append(TextEntity(m.group(0), "EMAIL", 0.99, "regex"))

    if NOR_PHONE_RE.fullmatch(t) and is_valid_phone(t):
        entities.append(TextEntity(t, "PHONE", 0.97, "rule"))
    elif expected_label == "PHONE":
        for m in GENERIC_PHONE_RE.finditer(t):
            if is_valid_phone(m.group(0)):
                entities.append(TextEntity(m.group(0), "PHONE", 0.90, "regex+context"))

    if POSTAL_RE.fullmatch(t):
        entities.append(TextEntity(t, "POSTAL_CODE", 0.96, "rule"))

    if looks_like_date(t):
        entities.append(TextEntity(t, "DATE", 0.93, "rule"))

    # ... additional rules omitted ...

    return dedupe_entities(entities)
\end{lstlisting}

Listing~\ref{lst:pii_detector_extract} shows that the final text-analysis
component combined multiple sources of evidence rather than relying only on the
NER model. Listing~\ref{lst:pii_detector_rules} illustrates how structured
entity types such as email addresses, phone numbers, postal codes, and dates
were reinforced through explicit rules in addition to model-based predictions.